% Actual class: 02
% Slide decks: M1-L3 DLL-Flow Control (2026); M1-L4 DLL-Error Control (2026)
% Exact scope: L3 pp. 1--32; L4 pp. 1--28
% NTULearn supplement: classroom recording transcript, used only to verify emphasis
% Human verified: Yes

\section{第 02 节课：Data Link Layer 的 Flow Control 与 Error Control}
\label{lecture:SC2008-L02}

\lecturemeta
  {SC2008-L02}
  {2026-08-18}
  {M1-L3 DLL-Flow Control (2026)；M1-L4 DLL-Error Control (2026)}
  {L3 pp.~1--32；L4 pp.~1--28}
  {课堂录像字幕（仅用于核对教师强调；原文未入库）}
  {已确认（2026-08-18）}

\subsection{本讲主线}

Flow Control（流量控制）防止 sender（发送方）过快发送而造成 receiver buffer overflow（接收缓冲区溢出）；Error Control（差错控制）处理 lost/damaged frame（丢失/损坏帧）和 lost ACK（确认丢失）。前者先假设传输无差错，后者通过 sequence number（序列号）、ACK/NAK 与 timeout（超时）把可靠性加入同一套机制。

\lecturetopic{SC2008-L02-T01}{ACK 与 Data Link Layer 的控制目标}

ACK（Acknowledgement，确认）是 receiver 根据接收状态新生成的 link-layer control message（链路层控制消息），不是 data frame（数据帧）解封装后剩余的内容。独立 ACK 通常具有自己的 link-layer framing（链路层成帧）和 error-detection field（差错检测字段），但一般不需要 Network Layer header（网络层首部）；也可 piggyback（捎带）在反向 data frame 中。

\begin{figure}[htbp]
  \centering
  \resizebox{0.92\textwidth}{!}{\begin{tikzpicture}[
  >=Latex,
  box/.style={draw=noteBlue,rounded corners=2pt,fill=blue!5,
    minimum height=1.05cm,align=center,text width=3.4cm},
  smallbox/.style={draw=noteBlue,rounded corners=2pt,fill=white,
    minimum height=0.9cm,align=center,text width=2.8cm},
  flow/.style={->,thick,draw=noteBlue}
]
  \node[smallbox] (frame) at (0,0)
    {$H_L\mid$ IP datagram $\mid T_L$\\data frame（数据帧）};
  \node[box] (dll) at (4.9,0)
    {Receiver Data Link Layer\\检查 FCS 与 sequence number};
  \node[smallbox] (network) at (9.7,1.5)
    {IP datagram\\交给 Network Layer};
  \node[smallbox] (ack) at (9.7,-1.5)
    {$ACK_n$ control frame\\下一步期待 $F_n$};

  \draw[flow] (frame) -- node[above]{接收并解析} (dll);
  \draw[flow] (dll) -- node[above,sloped,font=\small]{解封装} (network);
  \draw[flow] (dll) -- node[below,sloped,font=\small]{生成 ACK} (ack);
\end{tikzpicture}
}
  \caption{Receiver 对 data frame 解封装，并另外生成 ACK control frame；两者不是同一份数据。}
  \label{fig:sc2008-l02-ack-generation}
\end{figure}

本讲用 $ACK_n$ 表示 receiver 已累计确认当前编号范围内截至 $F_{n-1}$ 的 frames，下一步期待 $F_n$。例如 $ACK_3$ 表示 $F_0,F_1,F_2$ 已确认，下一帧期待 $F_3$。

\lecturetopic{SC2008-L02-T02}{Stop-and-Wait 与 Sliding Window Flow Control}

Stop-and-Wait（停止等待）每次只允许一个 outstanding frame（已发送但未确认的帧）。设
\[
  T_{\mathrm{frame}}=\frac{L}{R},\qquad
  T_{\mathrm{prop}}=\frac{d}{v},\qquad
  a=\frac{T_{\mathrm{prop}}}{T_{\mathrm{frame}}},
\]
其中 $L$ 为 frame length（帧长），$R$ 为 transmission rate（传输速率），$d$ 为距离，$v$ 为 signal propagation velocity（信号传播速度）。在 saturated input（持续有数据可发）、无差错并忽略 ACK transmission time 与 processing time 的条件下，
\[
  \boxed{U_{\mathrm{Stop-and-Wait}}=\frac{1}{1+2a}},
  \qquad
  \text{Throughput}=R\,U.
\]

Sliding Window（滑动窗口）允许最多 $N$ 个 outstanding frames 同时存在。Sequence number 使用 $k$ bits 并按 modulo $2^k$（模 $2^k$）循环；在本讲的无差错模型中 $N\leq 2^k$。其利用率为
\[
  \boxed{U_{\mathrm{Sliding\ Window}}
    =\min\!\left(1,\frac{N}{1+2a}\right)}.
\]
当 $N\geq 1+2a$ 时，sender 在首个 ACK 返回前仍有 frame 可发，链路可连续工作。

\begin{figure}[htbp]
  \centering
  \resizebox{0.96\textwidth}{!}{\begin{tikzpicture}[
  >=Latex,
  lifeline/.style={draw=black!55,thick},
  data/.style={->,very thick,draw=noteBlue},
  ack/.style={->,thick,draw=ntuRed},
  labelbox/.style={align=center,font=\small}
]
  \node[labelbox,font=\bfseries] at (2.5,0.75) {Stop-and-Wait};
  \node[labelbox] at (0,0) {Sender};
  \node[labelbox] at (5,0) {Receiver};
  \draw[lifeline] (0,-0.25) -- (0,-5.2);
  \draw[lifeline] (5,-0.25) -- (5,-5.2);
  \draw[data] (0,-0.65) -- (5,-1.65);
  \node[font=\scriptsize,anchor=south west,text=noteBlue] at (0.35,-0.65) {$F_0$};
  \draw[ack] (5,-2.05) -- (0,-3.05);
  \node[font=\scriptsize,text=ntuRed] at (2.5,-2.35) {$ACK_1$};
  \draw[data] (0,-3.55) -- (5,-4.55);
  \node[font=\scriptsize,anchor=south west,text=noteBlue] at (0.35,-3.55) {$F_1$};
  \draw[<->,draw=black!65] (-0.65,-0.5) -- node[left,align=right]
    {$T_{\mathrm{frame}}$\\$+2T_{\mathrm{prop}}$} (-0.65,-3.05);

  \node[labelbox,font=\bfseries] at (11.5,0.75) {Sliding Window（$N\ge 3$）};
  \node[labelbox] at (8,0) {Sender};
  \node[labelbox] at (15,0) {Receiver};
  \draw[lifeline] (8,-0.25) -- (8,-5.2);
  \draw[lifeline] (15,-0.25) -- (15,-5.2);
  \draw[data] (8,-0.65) -- (15,-1.65);
  \node[font=\scriptsize,anchor=south west,text=noteBlue] at (8.35,-0.65) {$F_0$};
  \draw[data] (8,-1.25) -- (15,-2.25);
  \node[font=\scriptsize,anchor=south west,text=noteBlue] at (8.35,-1.25) {$F_1$};
  \draw[data] (8,-1.85) -- (15,-2.85);
  \node[font=\scriptsize,anchor=south west,text=noteBlue] at (8.35,-1.85) {$F_2$};
  \draw[ack] (15,-2.15) -- (8,-3.15);
  \node[font=\scriptsize,text=ntuRed,fill=white,inner sep=1pt] at (11.5,-2.65) {$ACK_1$};
  \draw[data] (8,-3.65) -- (15,-4.65);
  \node[font=\scriptsize,anchor=south west,text=noteBlue] at (8.35,-3.65) {$F_3$};
  \node[labelbox,text=noteBlue] at (11.5,-5.45)
    {等待 ACK 时仍可发送，窗口足够大则没有 idle gap（空闲间隔）};
\end{tikzpicture}
}
  \caption{Stop-and-Wait 与 Sliding Window 的时序差异。斜线表示 propagation delay（传播时延），不是额外的 frame transmission time。}
  \label{fig:sc2008-l02-flow-timing}
\end{figure}

Receiver 可用 RNR（Receive Not Ready，接收方未准备好）暂时停止新传输；piggybacking（捎带确认）则把 ACK 放入反向 data frame，减少独立 ACK overhead（开销）。

\textbf{对应例题：}\relatedexercise{SC2008-L02-E01}{Stop-and-Wait link utilization}。

\lecturetopic{SC2008-L02-T03}{Error Detection、FEC 与 ARQ}

Lost frame（丢失帧）表示 receiver 未获得可识别的 frame；damaged frame（损坏帧）表示 frame 到达但部分 bits 出错。Parity check（奇偶校验）能检测任意奇数个 bit flips（比特翻转），但可能漏掉偶数个；CRC（Cyclic Redundancy Check，循环冗余校验）主要用于 error detection（差错检测），不能直接等同于 error correction（纠错）。本讲将这部分标为不可考。

FEC（Forward Error Correction，前向纠错）用冗余信息直接纠正一定范围内的错误；ARQ（Automatic Repeat Request，自动重传请求）则在检测到失败后请求重传。后续可考内容集中在 Stop-and-Wait、Go-Back-N 与 Selective Reject/Repeat 三种 ARQ。

\lecturetopic{SC2008-L02-T04}{Stop-and-Wait ARQ}

Stop-and-Wait ARQ 使用交替的 sequence numbers $0,1$：frame 丢失时 sender timeout 后重传；damaged frame 可被丢弃并触发 NAK；ACK 丢失时 sender 会重传相同 frame，receiver 利用 sequence number 识别 duplicate（重复帧），丢弃重复数据并再次返回 ACK。

若 $P$ 表示一次 transmission attempt（传输尝试）未产生有效交付的概率，则
\[
  \boxed{U_{\mathrm{Stop-and-Wait\ ARQ}}
    =\frac{1-P}{1+2a}}.
\]
因 retransmission（重传）本身也可能再次失败，每成功交付一个 frame 的期望尝试次数为
\[
  1+P+P^2+\cdots=\frac{1}{1-P},
\]
所以有效 attempts 的长期比例是 $1-P$，而不是 $1/(1+P)$。

\textbf{对应例题：}\relatedexercise{SC2008-L02-E02}{Stop-and-Wait ARQ with frame loss}。

\lecturetopic{SC2008-L02-T05}{Go-Back-N 与 Selective Reject ARQ}

\begin{center}
\small
\begin{tabularx}{\textwidth}{@{}>{\raggedright\arraybackslash}p{3.4cm}XX@{}}
  \toprule
  特征 & Go-Back-N（GBN） & Selective Reject/Repeat（SR） \\
  \midrule
  Out-of-order frame & 丢弃 & 缓存 \\
  出错后的重传 & 错误帧及全部后续帧 & 只重传错误帧 \\
  Receiver complexity & 较低 & 较高 \\
  $k$-bit 最大窗口 & $2^k-1$ & $2^{k-1}$ \\
  \bottomrule
\end{tabularx}
\end{center}

若 $F_1$ 丢失而 $F_2,F_3$ 到达，GBN receiver 丢弃 $F_2,F_3$，sender 重传 $F_1,F_2,F_3$；SR receiver 缓存 $F_2,F_3$，sender 只重传 $F_1$，随后 receiver 重排并交付 $F_1,F_2,F_3$。

Window-size restriction（窗口限制）防止旧 frame 与编号回绕后的新 frame 混淆。GBN receiver window size 为 1，因此
\[
  W_S+1\leq 2^k
  \quad\Longrightarrow\quad
  \boxed{W_S\leq 2^k-1}.
\]
当 SR 的 sender/receiver windows 都为 $N$ 时，需要旧 sender window 与下一 receiver window 不重叠：
\[
  2N\leq 2^k
  \quad\Longrightarrow\quad
  \boxed{N\leq 2^{k-1}}.
\]

\begin{figure}[htbp]
  \centering
  \resizebox{0.98\textwidth}{!}{\begin{tikzpicture}[
  cell/.style={draw=black!65,minimum width=0.72cm,minimum height=0.62cm,
    inner sep=0pt,align=center},
  old/.style={cell,fill=blue!13},
  next/.style={cell,fill=green!16},
  overlap/.style={cell,fill=orange!35},
  rowlabel/.style={anchor=east,align=right,font=\small}
]
  \node[rowlabel] at (-0.3,0.65) {GBN, $k=3,N=7$\\old sender window};
  \foreach \x in {0,...,6} {\node[old] at (0.72*\x,0.65) {\x};}
  \node[next] at (5.04,0.65) {7};
  \node[anchor=west,font=\small,text=noteBlue] at (5.55,0.65)
    {$W_R$ 只期待 7，旧 $F_0$ 被拒绝};

  \node[rowlabel] at (-0.3,-0.65) {SR safe, $k=3,N=4$\\$W_S^{\mathrm{old}}$};
  \foreach \x in {0,...,3} {\node[old] at (0.72*\x,-0.65) {\x};}
  \node[rowlabel] at (5.25,-0.65) {$W_R^{\mathrm{next}}$};
  \node[next] at (5.65,-0.65) {4};
  \node[next] at (6.37,-0.65) {5};
  \node[next] at (7.09,-0.65) {6};
  \node[next] at (7.81,-0.65) {7};
  \node[anchor=west,font=\small,text=noteBlue] at (8.65,-0.65) {disjoint（不重叠）};

  \node[rowlabel] at (-0.3,-1.95) {SR unsafe, $k=3,N=5$\\$W_S^{\mathrm{old}}$};
  \foreach \x in {0,...,4} {\node[old] at (0.72*\x,-1.95) {\x};}
  \node[rowlabel] at (5.25,-1.95) {$W_R^{\mathrm{next}}$};
  \node[next] at (5.65,-1.95) {5};
  \node[next] at (6.37,-1.95) {6};
  \node[next] at (7.09,-1.95) {7};
  \node[overlap] at (7.81,-1.95) {0};
  \node[overlap] at (8.53,-1.95) {1};
  \node[anchor=west,font=\small,text=verifyOrange] at (9.05,-1.95)
    {旧 $F_0,F_1$ 可能被当作新 frame};
\end{tikzpicture}
}
  \caption{ACK 全部丢失时 sender 仍停留在旧窗口，而 receiver 已进入下一窗口；安全性取决于两组 sequence numbers 是否重叠。}
  \label{fig:sc2008-l02-window-ambiguity}
\end{figure}

重复丢失只延长等待时间：未收到 ACK 时 sender 不能进入第三个编号世代，receiver 也不会因窗口外的 duplicate frame 再次前移。

Selective Reject 的课程性能模型为
\[
  \boxed{U_{\mathrm{SR}}
    =(1-P)\min\!\left(1,\frac{N}{1+2a}\right)}.
\]
Go-Back-N 的精确公式不要求记忆；本课程要求相关数值计算使用 SR 公式。这是 course-specific simplification（课程特定简化），不代表两种协议的实际性能完全相同。

\subsection{一分钟回顾}

\[
\begin{aligned}
  U_{\mathrm{Stop-and-Wait}}&=\frac{1}{1+2a},
  &U_{\mathrm{Sliding\ Window}}&=\min\!\left(1,\frac{N}{1+2a}\right),\\
  U_{\mathrm{Stop-and-Wait\ ARQ}}&=\frac{1-P}{1+2a},
  &U_{\mathrm{SR}}&=(1-P)\min\!\left(1,\frac{N}{1+2a}\right).
\end{aligned}
\]
\[
  \boxed{N_{\max}^{\mathrm{Flow}}=2^k,\qquad
  N_{\max}^{\mathrm{GBN}}=2^k-1,\qquad
  N_{\max}^{\mathrm{SR}}=2^{k-1}}.
\]
